\chapter{Przestrzenie euklidesowe i unitarne}
\label{ch:lin12-euklidesowe-unitarne}
\index{przestrzeń euklidesowa@Przestrzeń euklidesowa}
\index{przestrzeń unitarna@Przestrzeń unitarna}
\index{forma hermitowska@Forma hermitowska}
\index{iloczyn skalarny zespolony@Iloczyn skalarny zespolony}

\noindent\textbf{Cel rozdziału.}
W Rozdziale~\ref{ch:lin9-formy-kwadratowe} udowodniliśmy
podstawowe twierdzenia o rzeczywistych iloczynach skalarnych:
nierówność Cauchy'ego--Schwarza, ortogonalizację
Grama--Schmidta i istnienie rzutów ortogonalnych.
Nie będziemy powtarzać tych dowodów.
Rozszerzymy natomiast teorię na przestrzenie zespolone.
Właściwym odpowiednikiem rzeczywistej symetrycznej formy
dwuliniowej jest tam \emph{forma hermitowska}, która
w pierwszym argumencie jest sprzężenie liniowa,
a w drugim liniowa. Zbadamy polaryzację zespoloną,
współrzędne, nierówność Cauchy'ego--Schwarza,
bazy ortonormalne, reprezentację Riesza i
odwzorowania unitarne. W kolejnych rozdziałach
będziemy mogli dzięki temu badać bardziej szczegółowo
rzuty i operatory na przestrzeniach z iloczynem skalarnym.

Przez \(K\) oznaczamy \(\mathbb R\) albo \(\mathbb C\).
Symbol \(\overline a\) oznacza sprzężenie zespolone,
a dla \(K=\mathbb R\) przyjmujemy \(\overline a=a\).
Nie wymagamy skończonego wymiaru, o ile założenie
to nie jest wyraźnie wskazane. Przestrzenie są
algebraiczne; nie wprowadzamy tu pojęcia zupełności
przestrzeni unormowanej.

\section{Formy hermitowskie i iloczyny skalarne}
\label{sec:lin12-hermitowskie}

\begin{definicja}[Forma półtoraliniowa i hermitowska]
\label{def:lin12-hermitowska}
Niech \(V\) będzie przestrzenią wektorową nad \(K\).
Odwzorowanie \(h\colon V\times V\to K\)
nazywamy \emph{formą półtoraliniową}
(z konwencją liniowości w drugim argumencie),
jeśli dla wszystkich \(u,v,w\in V\)
oraz \(a,b\in K\) zachodzą
\[
 \begin{aligned}
 h(au+bv,w)&=\overline a\,h(u,w)+
              \overline b\,h(v,w),\\
 h(w,au+bv)&=a\,h(w,u)+b\,h(w,v).
 \end{aligned}
\]
Formę taką nazywamy \emph{hermitowską}, jeśli
\[
 h(v,u)=\overline{h(u,v)}
 \qquad(u,v\in V).
\]
Forma hermitowska jest \emph{dodatnio określona},
jeżeli \(h(v,v)>0\) dla każdego \(v\ne0_V\).
Wówczas nazywamy ją \emph{iloczynem skalarnym},
a parę \((V,h)\) \emph{przestrzenią euklidesową}
dla \(K=\mathbb R\) albo \emph{przestrzenią unitarną}
dla \(K=\mathbb C\).
\end{definicja}

\begin{uwaga}[Konwencja dotycząca argumentów]
\label{rem:lin12-konwencja}
W literaturze występuje też konwencja, w której
iloczyn skalarny jest liniowy w pierwszym,
a sprzężenie liniowy w drugim argumencie.
Nie wolno mieszać obu konwencji:
przy naszej definicji
\(h(u,av)=a\,h(u,v)\),
natomiast \(h(au,v)=\overline a\,h(u,v)\).
Wszystkie wzory na macierze, współrzędne,
polaryzację oraz reprezentację Riesza
są poniżej zgodne właśnie z tą konwencją.
\end{uwaga}

\begin{stwierdzenie}[Przekątna formy hermitowskiej]
\label{prop:lin12-przekatna}
Niech \(h\) będzie formą hermitowską na \(V\).
Wtedy \(h(v,v)\in\mathbb R\) dla każdego \(v\in V\).
Ponadto \(h(0_V,v)=h(v,0_V)=0_K\).
Jeżeli \(h\) jest dodatnio określona, to
\(h(v,v)=0\) zachodzi wyłącznie dla \(v=0_V\).
\end{stwierdzenie}
\begin{proof}
Z hermitowskości dla \(u=v\) otrzymujemy
\(h(v,v)=\overline{h(v,v)}\), więc
\(h(v,v)\) jest rzeczywista.
Półtoraliniowość daje
\(h(0_V,v)=0_K\) i \(h(v,0_V)=0_K\)
po podstawieniu skalaru \(0_K\).
Pozostała teza wynika bezpośrednio
z dodatniej określoności.
\end{proof}

\begin{przyklad}[Standardowy iloczyn skalarny]
\label{prz:lin12-standardowy}
Niech \(n\in\mathbb N_0\).
Na \(K^n\) określamy
\[
 h(z,w)=\sum_{j=1}^{n}\overline{z_j}w_j.
\]
Pierwszy argument jest sprzężenie liniowy,
a drugi liniowy. Zamiana argumentów
daje sprzężenie wartości całej sumy.
Ponadto
\(h(z,z)=\sum_j|z_j|^2\ge0\),
przy czym równość zachodzi dokładnie dla \(z=0\).
Otrzymujemy przestrzeń euklidesową
dla \(K=\mathbb R\) i unitarną
dla \(K=\mathbb C\).
Przy \(n=0\) jest to jedyny iloczyn
skalarny na przestrzeni zerowej;
warunek dodatniej określoności jest
spełniony w sposób pusty.
\end{przyklad}

\begin{przyklad}[Dlaczego zwykła dwuliniowość nie wystarcza]
\label{prz:lin12-nie-dwuliniowa}
Na \(\mathbb C\) funkcja \(b(z,w)=zw\)
jest dwuliniowa nad \(\mathbb C\),
ale nie jest dodatnio określonym
iloczynem skalarnym:
\(b(i,i)=-1\).
Ponadto \(b(1,i)=i=b(i,1)\),
podczas gdy hermitowskość wymagałaby
\(b(i,1)=\overline{b(1,i)}=-i\).
Poprawny iloczyn skalarny na \(\mathbb C\)
ma postać \(h(z,w)=\overline zw\).
\end{przyklad}

\section{Polaryzacja zespolona}
\label{sec:lin12-polaryzacja}

\begin{twierdzenie}[Tożsamość polaryzacyjna]
\label{thm:lin12-polaryzacja}
Niech \(h\colon V\times V\to\mathbb C\)
będzie formą hermitowską i
\(q(v):=h(v,v)\in\mathbb R\).
Wówczas dla wszystkich \(u,v\in V\)
\[
 \begin{aligned}
 h(u,v)=\frac14\bigl(
 &q(u+v)-q(u-v)\\
 &-i\,q(u+iv)+i\,q(u-iv)\bigr).
 \end{aligned}
\]
W szczególności dwie formy hermitowskie
na tej samej przestrzeni zespolonej,
które mają jednakowe wartości na przekątnej,
są równe. Dla \(K=\mathbb R\) odpowiedni wzór
ma postać
\[
 h(u,v)=\tfrac14\bigl(q(u+v)-q(u-v)\bigr).
\]
Nie trzeba zakładać dodatniej określoności.
\end{twierdzenie}
\begin{proof}
Niech \(z=h(u,v)\).
Z półtoraliniowości i hermitowskości
otrzymujemy
\[
 \begin{aligned}
 q(u+v)&=q(u)+q(v)+z+\overline z,\\
 q(u-v)&=q(u)+q(v)-z-\overline z.
 \end{aligned}
\]
Zatem różnica pierwszych dwóch wartości
wynosi \(2(z+\overline z)=4\operatorname{Re}z\).
Dla argumentów \(u\pm iv\) mamy
\[
 \begin{aligned}
 q(u+iv)&=q(u)+q(v)+iz-i\overline z,\\
 q(u-iv)&=q(u)+q(v)-iz+i\overline z.
 \end{aligned}
\]
Ich różnica jest równa
\(2i(z-\overline z)=-4\operatorname{Im}z\).
Wstawiając te dwie różnice do wzoru z tezy,
otrzymujemy
\(\operatorname{Re}z+i\operatorname{Im}z=z\).

Jeśli \(h_1(u,u)=h_2(u,u)\)
dla wszystkich \(u\), to zastosowanie
tego samego wzoru do obu form
daje \(h_1(u,v)=h_2(u,v)\)
dla wszystkich \(u,v\).
Dla \(K=\mathbb R\)
formy hermitowskie są symetrycznymi
formami dwuliniowymi, więc wystarczają
pierwsze dwie równości.
\end{proof}

\begin{przyklad}[Odzyskanie części urojonej]
\label{prz:lin12-polaryzacja-i}
Niech \(V=\mathbb C\) z iloczynem
\(h(z,w)=\overline zw\), a
\(q(z)=|z|^2\).
Dla \(u=1\), \(v=i\)
mamy \(h(1,i)=i\), lecz
\(q(1+i)=q(1-i)=2\).
Rzeczywista różnica kwadratów jest zerowa.
Natomiast
\(q(1+i\cdot i)=q(0)=0\)
i \(q(1-i\cdot i)=q(2)=4\).
Tożsamość polaryzacyjna daje
\(\frac14(0-i(0-4))=i\),
zgodnie z wartością iloczynu skalarnego.
Zastosowanie wyłącznie rzeczywistego
wzoru polaryzacyjnego prowadziłoby
do błędnego wyniku \(0\).
\end{przyklad}

\section{Macierz Grama i zmiana bazy}
\label{sec:lin12-gram}

\begin{definicja}[Macierz Grama]
\label{def:lin12-gram}
Niech \(\dim V=n<\infty\), a
\(B=(e_1,\ldots,e_n)\) będzie
uporządkowaną bazą \(V\).
Dla formy hermitowskiej \(h\)
\emph{macierzą Grama} nazywamy
\[
 G_B(h)=(h(e_i,e_j))_{i,j=1}^{n}.
\]
Dla macierzy \(A\) symbol \(A^*\)
oznacza macierz sprzężoną i transponowaną,
tj. \((A^*)_{ij}=\overline{A_{ji}}\).
Macierz \(A\) jest \emph{hermitowska},
jeżeli \(A^*=A\).
\end{definicja}

\begin{twierdzenie}[Współrzędne i kongruencja hermitowska]
\label{thm:lin12-gram-zmiana}
Niech \(V\) będzie skończenie wymiarową
przestrzenią nad \(K\),
\(h\) formą hermitowską oraz
\(B=(e_1,\ldots,e_n)\) bazą.
Wtedy macierz \(G=G_B(h)\)
jest hermitowska, a dla
\(X=[u]_B\), \(Y=[v]_B\) zachodzi
\[
 h(u,v)=X^*GY.
\]
Odwrotnie, każda macierz hermitowska
\(G\in\operatorname{Mat}_{n\times n}(K)\)
definiuje w ten sposób jednoznaczną
formę hermitowską na \(K^n\).
Jeżeli baza \(B'\) jest wyrażona
w bazie \(B\) macierzą przejścia \(P\),
czyli \([v]_B=P[v]_{B'}\), to
\[
 G_{B'}(h)=P^*G_B(h)P.
\]
Forma \(h\) jest dodatnio określona
dokładnie wtedy, gdy
\(X^*GX>0\) dla wszystkich \(X\ne0\).
\end{twierdzenie}
\begin{proof}
Dla \(G_{ij}=h(e_i,e_j)\)
hermitowskość daje
\(G_{ji}=\overline{G_{ij}}\),
co jest równoważne \(G^*=G\).
Pisząc
\(u=\sum_i x_ie_i\), \(v=\sum_j y_je_j\),
stosujemy półtoraliniowość:
\[
 h(u,v)=\sum_{i,j}\overline{x_i}\,
         h(e_i,e_j)y_j=X^*GY.
\]
Jeśli \(G^*=G\), to
\(h_G(X,Y)=X^*GY\)
jest sprzężenie liniowe w \(X\)
i liniowe w \(Y\), a ponadto
\[
 \overline{h_G(Y,X)}
 =\overline{Y^*GX}
 =X^*G^*Y=X^*GY.
\]
Zatem \(h_G\) jest hermitowska.
Jednoznaczność wynika z jej wartości
na parach wektorów bazowych.
W nowej bazie
\(X=P X'\), \(Y=P Y'\),
więc
\(X^*GY=(X')^*P^*GPY'\);
to daje wzór kongruencji.
Ostatnia równoważność jest dokładnie
definicją dodatniej określoności
w nieosobliwym układzie współrzędnych.
\end{proof}

\begin{przyklad}[Niediagonalna macierz hermitowska]
\label{prz:lin12-gram-i}
Niech na \(\mathbb C^2\) zadana będzie
macierz
\[
 G=\begin{pmatrix}2&i\\-i&2\end{pmatrix}.
\]
Jest hermitowska, bo \(G^*=G\),
ale nie jest macierzą rzeczywistą.
Odpowiadająca jej forma ma
na przekątnej wartość
\[
 q(z)=2|z_1|^2+2|z_2|^2
       -2\operatorname{Im}(\overline{z_1}z_2).
\]
Po uzupełnieniu kwadratu
\[
 q(z)=2\left|z_1+\frac i2z_2\right|^2
       +\frac32|z_2|^2.
\]
Oba składniki są nieujemne,
a ich suma znika wyłącznie
dla \(z_1=z_2=0\).
Forma jest zatem dodatnio określona.
Z tego samego rozkładu wynika,
dlaczego hermitowskość \(G\)
nie wystarcza sama w sobie
do dodatniej określoności:
np. macierz \(\operatorname{diag}(1,-1)\)
jest hermitowska, lecz nie jest
dodatnio określona.
\end{przyklad}

\section{Nierówność Cauchy'ego--Schwarza i norma}
\label{sec:lin12-norma}

\begin{twierdzenie}[Nierówność Cauchy'ego--Schwarza nad \(\mathbb C\)]
\label{thm:lin12-cauchy-schwarz}
Niech \((V,h)\) będzie przestrzenią unitarną.
Dla dowolnych \(u,v\in V\) mamy
\[
 |h(u,v)|^2\le h(u,u)\,h(v,v).
\]
Równość zachodzi dokładnie wtedy,
gdy \(u,v\) są liniowo zależne nad \(\mathbb C\).
Funkcja
\(\|v\|:=\sqrt{h(v,v)}\)
jest normą na \(V\), a
\[
 \|u+v\|\le\|u\|+\|v\|.
\]
\end{twierdzenie}
\begin{proof}
Jeżeli \(v=0_V\), nierówność
jest równością i oba wektory
są liniowo zależne.
Niech więc \(v\ne0_V\).
Liczba \(h(v,v)\) jest dodatnia.
Kładziemy
\(c:=h(v,u)/h(v,v)\).
Wówczas
\(h(v,u-cv)=h(v,u)-c\,h(v,v)=0\).
Z nieujemności kwadratu normy
otrzymujemy
\[
 \begin{aligned}
 0&\le h(u-cv,u-cv)\\
  &=h(u,u)-\frac{|h(u,v)|^2}{h(v,v)}.
 \end{aligned}
\]
Pomnożenie przez \(h(v,v)>0\)
daje żądaną nierówność.
Równość jest równoważna
\(h(u-cv,u-cv)=0\),
czyli \(u=cv\).
W przypadku \(v=0\)
warunek zależności również zachodzi.

Dodatnia określoność daje
nieujemność \(\|v\|\) i
\(\|v\|=0\iff v=0\).
Półtoraliniowość zapewnia
\[
 \|av\|^2=h(av,av)=|a|^2 h(v,v),
\]
więc \(\|av\|=|a|\|v\|\).
Wreszcie
\[
 \begin{aligned}
 \|u+v\|^2
 &=\|u\|^2+2\operatorname{Re}h(u,v)+\|v\|^2\\
 &\le \|u\|^2+2|h(u,v)|+\|v\|^2\\
 &\le (\|u\|+\|v\|)^2.
 \end{aligned}
\]
Obie strony są nieujemne,
więc pierwiastkowanie daje
nierówność trójkąta.
\end{proof}

\begin{wniosek}[Wzór Pitagorasa i równoległoboku]
\label{cor:lin12-pitagoras}
Niech \((V,h)\) będzie przestrzenią
euklidesową lub unitarną.
Jeśli \(h(u,v)=0\), to
\[
 \|u+v\|^2=\|u\|^2+\|v\|^2.
\]
Dla dowolnych \(u,v\in V\) zachodzi
\[
 \|u+v\|^2+\|u-v\|^2
 =2\|u\|^2+2\|v\|^2.
\]
\end{wniosek}
\begin{proof}
Rozwijamy kwadraty norm za pomocą
półtoraliniowości. W pierwszym
przypadku oba mieszane składniki
\(h(u,v)\) i \(h(v,u)\) znikają.
W drugim przypadku składniki
mieszane w dwóch rozwinięciach
są przeciwne i znoszą się po dodaniu.
\end{proof}

\section{Bazy ortonormalne w przestrzeniach unitarnych}
\label{sec:lin12-bazy}

\begin{definicja}[Rodzina ortonormalna]
\label{def:lin12-ortonormalna}
W przestrzeni \((V,h)\) rodzinę
\((e_j)_{j\in J}\) nazywamy
\emph{ortonormalną}, jeśli
\(h(e_i,e_j)=\delta_{ij}\)
dla wszystkich \(i,j\in J\).
Jeśli jest bazą Hamela,
nazywamy ją \emph{bazą ortonormalną}.
\end{definicja}

\begin{twierdzenie}[Ortogonalizacja Grama--Schmidta nad \(\mathbb C\)]
\label{thm:lin12-gram-schmidt}
Niech \((V,h)\) będzie przestrzenią
unitarną, a \(v_1,\ldots,v_m\in V\),
\(m\ge0\), wektorami liniowo niezależnymi.
Istnieje rodzina ortonormalna
\(e_1,\ldots,e_m\), dla której
\[
 \operatorname{span}\{e_1,\ldots,e_j\}
 =\operatorname{span}\{v_1,\ldots,v_j\}
 \qquad(1\le j\le m).
\]
Można ją skonstruować wzorami
\[
 w_j=v_j-\sum_{\ell=1}^{j-1}
            h(e_\ell,v_j)e_\ell,
 \qquad
 e_j=\frac{w_j}{\sqrt{h(w_j,w_j)}}.
\]
W szczególności każda skończenie
wymiarowa przestrzeń unitarna
posiada bazę ortonormalną.
\end{twierdzenie}
\begin{proof}
Dla \(m=0\) rodzina pusta
spełnia tezę. Dla \(m\ge1\)
przeprowadzamy indukcję.
Przy \(j=1\) wektor \(v_1\)
jest niezerowy i normalizujemy go.
Załóżmy, że
\(e_1,\ldots,e_{j-1}\)
są ortonormalne i
rozpinają tę samą przestrzeń,
co \(v_1,\ldots,v_{j-1}\).
Dla każdego \(i<j\)
półtoraliniowość w drugim
argumencie daje
\[
 \begin{aligned}
 h(e_i,w_j)
 &=h(e_i,v_j)
   -\sum_{\ell<j}h(e_\ell,v_j)h(e_i,e_\ell)\\
 &=h(e_i,v_j)-h(e_i,v_j)=0.
 \end{aligned}
\]
Jeśli \(w_j=0\), to \(v_j\)
jest kombinacją poprzednich \(e_i\),
a zatem i poprzednich \(v_i\),
wbrew niezależności.
Dlatego \(h(w_j,w_j)>0\)
i wzór na \(e_j\) ma sens.
Nowy wektor jest normy \(1\)
i ortogonalny do poprzednich.
Ponadto \(w_j-v_j\) należy
do poprzedniej rozpiętości,
więc obie rodziny mają
taką samą rozpiętość do kroku \(j\).
Indukcja kończy konstrukcję.
Przy skończonym wymiarze
stosujemy ją do dowolnej bazy.
\end{proof}

\begin{przyklad}[Ortonormalizacja dwóch wektorów zespolonych]
\label{prz:lin12-gs-c2}
W \(\mathbb C^2\) ze standardowym
iloczynem skalarnym rozważmy
\(v_1=(1,i)\), \(v_2=(1,0)\).
Pierwszy wektor ortonormalny to
\[
 e_1=\frac1{\sqrt2}(1,i).
\]
Ponieważ
\(h(e_1,v_2)=1/\sqrt2\),
otrzymujemy
\[
 w_2=v_2-\tfrac1{\sqrt2}e_1
    =\bigl(\tfrac12,-\tfrac i2\bigr),
 \qquad
 e_2=\tfrac1{\sqrt2}(1,-i).
\]
Sprawdzenie bezpośrednie daje
\(h(e_1,e_2)=(1+(-i)(-i))/2=0\)
i \(h(e_2,e_2)=1\).
W tym rachunku sprzężenie
w pierwszym argumencie
jest istotne: zwykły iloczyn
współrzędnych nie dałby zera.
\end{przyklad}
